\documentclass[11pt,a4paper]{article}
\usepackage[T1]{fontenc}
\usepackage{lmodern,amsmath,amssymb,booktabs,microtype}
\usepackage[margin=24mm]{geometry}
\usepackage[colorlinks=true,urlcolor=blue,linkcolor=blue]{hyperref}
\newcommand{\EE}{\mathbb E}
\newcommand{\cc}{\mathcal C}
\newcommand{\rr}{\mathcal R}
\newcommand{\core}{\mathcal E}
\newcommand{\JJ}{\mathcal J}
\DeclareMathOperator{\tr}{tr}
\title{Quadratic contractions of the memory core\\An exact adjacent-charge covariance formula}
\author{Alan Day \and Chengheng\thanks{AI research assistant (OpenAI).}}
\date{Local research note --- 6 November 2026}
\begin{document}\maketitle
\begin{abstract}
We evaluate the ordinary second Gaussian-chaos part of the quartic
operator entering the second centered memory coefficient. The answer
uses neutral and adjacent-charge embedded-ensemble covariance maps,
retaining the same couplings on both charge sectors. A finite example
has a nonzero dynamical residual after removing sample-dependent
scalar parts and the previously available quadratic descendants.
The formula supplies the missing finite-$N$ contraction interface;
its double-scaled residual strength is the next calculation.
\end{abstract}

\section{Model, target and the correct Gaussian subtraction}
Use standard CAR and the full overlapping complex Gaussian $p$-body
Hamiltonian. The couplings obey
$\EE J_{IJ}J_{KL}=w\delta_{IL}\delta_{JK}$. All
results here are exact at finite $N$, fixed charge $m\ge p$, with
annealed normalized Hilbert--Schmidt inner product
$\langle X,Y\rangle_m=\EE\tr_m(X^\dagger Y)$. Take $\hbar=1$.
The physical variance is
$w=4^p(J_s^2N/p^2)(p!/N^p)^2$; explicit finite tests use $w=1$.
No large-$N$ or small-$\lambda$ approximation is used below.

Write $\mathcal L=[H,\cdot]$ and
\begin{equation}
 A=\sum_i(\mathcal Lc_i)^\dagger\mathcal Lc_i,\qquad
 \core=\sum_i(\mathcal L^2c_i)^\dagger\mathcal L^2c_i
                  +\tfrac14\mathcal L^2A.
\end{equation}
The preceding memory calculation identified $\core$ as the core of
$mZ_2$ after its explicit first-return counterterm, up to a vanishing
double-scaled remainder. This note evaluates its lower coupling chaos.

Choose independent real Gaussian coordinates $x_a$ with variances
$v_a$, so that $H=\sum_a x_a V_a$. For any homogeneous quartic
matrix polynomial $P$, put $\Delta=\sum_a v_a\partial_{x_a}^2$ and
\begin{equation}
 Q_P=\tfrac12\Delta P,\qquad e_P=\EE P.
 \label{eq:heat}
\end{equation}
The Gaussian normal-ordering identity gives exactly
\begin{equation}
 W_0P=e_P,\qquad W_2P=Q_P-2e_P,\qquad W_4P=P-Q_P+e_P.
 \label{eq:chaos}
\end{equation}
Indeed $\Delta^2P=8e_P$ and $\EE Q_P=2e_P$.
These are coordinate-chaos projections. The $q$-Hermite Wick
polynomials of the spectral law in \cite{TT} have a different
definition; an equality between them requires a separate limit argument.

\newpage
\section{Charge removal and the covariance primitives}
On neighboring sectors write $h=H_m$, $k=H_{m-1}$ and
$c_i:\mathcal H_m\to\mathcal H_{m-1}$. Define
\begin{equation}
 \rr(X)=\sum_i c_i^\dagger Xc_i,\quad R_j=\rr(k^j),\quad
 R_0=mI,\quad R_1=(m-p)h.
\end{equation}
Directly expanding the two fermion derivatives gives
\begin{align}
 A&=R_2+(2p-m)h^2,\\
 \core&=R_4-2\{h,R_3\}+\tfrac54\{h^2,R_2\}
                  +\tfrac72 hR_2h+(4p-3m)h^4. \label{eq:raw}
\end{align}
For example the square of
$k^2c_i-2kc_ih+c_ih^2$ yields the coefficients
$1,-2,1,4,4p-3m$ before adding $\mathcal L^2A/4$.

Let $(u,v)$ be an independent copy of $(h,k)$, using the same
independent couplings on the two sectors. Set
\begin{align}
 \Phi_+(X)&=\EE_u[uXu],&\Phi_-(X)&=\EE_v[vXv],\\
 \Psi(X)&=\EE_{u,v}[uXv],&
 \Gamma(X,Y)&=\sum_i\Psi(c_i^\dagger X)Yc_i. \label{eq:maps}
\end{align}
Here $\Psi$ acts on $d_m\times d_{m-1}$ matrices, while $X,Y$
in $\Gamma$ act on $\mathcal H_{m-1}$. The neutral constants are
\begin{equation}
 C=w\binom mp\binom{N-m+p}p,\quad
 C_-=w\binom{m-1}p\binom{N-m+p+1}p,\quad
 a=(2p-m)C+mC_-.
\end{equation}
Thus $\Phi_+(I)=CI$, $\Phi_-(I)=C_-I$, $\EE A=aI$.
All binomial coefficients outside their range vanish.

These are the standard embedded-ensemble covariance operations
\cite{Kota,TT}, not new independent random probes. If
$T^{(r)}_{IJ}=(a_I^\dagger a_J)|_{\mathcal H_r}$
denotes the charge-$r$ block, the mixed map is
\begin{equation}
 \Psi(X)=w\sum_{I,J} T^{(m)}_{IJ}\,X\,T^{(m-1)}_{JI}.
 \label{eq:mixed}
\end{equation}
Neutral maps are obtained by using the same block on both sides.
In a neutral irreducible rank-$\nu$ channel,
\begin{equation}
 \Phi_m(O_\nu)=\kappa_\nu(m)O_\nu,\qquad
 \kappa_\nu(m)=w\binom{m-\nu}p\binom{N-m+p-\nu}p.
\end{equation}
The cross-sector map in \eqref{eq:mixed} must be retained before
taking an asymptotic limit.

\newpage
\section{The closed single-contraction formula}
For a polynomial $P(h,k)$ let
$\cc P=\EE_{u,v}[t^2]P(h+tu,k+tv)$. This equals $Q_P$ in
\eqref{eq:heat}. Introduce
\begin{align}
 P_+&=\Phi_+(h),&P_-&=\Phi_-(k),\\
 U&=\Gamma(k,I)+\Gamma(I,k),\\
 V&=\Gamma(I,k^2)+\Gamma(k,k)+\Gamma(k^2,I),\\
 U_h&=\sum_i\big[\Psi(hc_i^\dagger k)c_i
                         +\Psi(hc_i^\dagger)kc_i\big].
\end{align}
Every quantity on the right contains only the original $h,k$ and
the deterministic covariance maps. The five contractions in
\eqref{eq:raw} are
\begin{align}
 D_4:=\cc R_4
   &=\rr\big(3C_-k^2+\{k,P_-\}+\Phi_-(k^2)\big),\\
 D_{13}:=\cc\{h,R_3\}
   &=\{h,\rr(2C_-k+P_-)\}+V+V^\dagger,\\
 D_{22}:=\cc\{h^2,R_2\}
   &=2CR_2+2mC_-h^2+hU+U_h+U^\dagger h+U_h^\dagger,\\
 D_{121}:=\cc(hR_2h)
   &=\Phi_+(R_2)+Uh+hU^\dagger+mC_-h^2,\\
 D_{40}:=\cc h^4
   &=3Ch^2+\{h,P_+\}+\Phi_+(h^2).
\end{align}
Consequently the answer is the Hermitian quadratic polynomial
\begin{equation}
 \boxed{Q_{\core}=D_4-2D_{13}+\tfrac54D_{22}
                  +\tfrac72D_{121}+(4p-3m)D_{40}.}
 \label{eq:answer}
\end{equation}
The scalar $e=\EE\core=e_mI$ follows either by a second contraction
or by the original fourth moment. Equation \eqref{eq:chaos} then
recovers $W_2\core$ and $W_4\core$ exactly.

To check the ordering, write
$R_2(h+tu,k+tv)=R_2+tB+t^2\rr(v^2)$ with
$B=\rr(kv+vk)$. Then $\EE uB=U$,
$\EE uhB=U_h$, and $\EE Bu=U^\dagger$. These three expectations
generate the mixed terms in $D_{22}$ and $D_{121}$. For $R_3$, the
coefficient linear in $v$ is $\rr(vk^2+kvk+k^2v)$;
its left contraction with $u$ is $V$. This proves the mixed
formulas without an independence assumption between neighboring sectors.

The formula is also an interface for restoring the low-chaos
parts of the four-slot descendant basis: all of its members are
linear combinations of the five raw words in \eqref{eq:raw}.

\newpage
\section{Projection onto old directions and a dynamical residual}
Put $B=h^2-CI$ and $O=A-aI$. The exact projection coefficient
onto the energy quadratic is
\begin{equation}
 c_E=\frac{\langle B,W_2\core\rangle_m}{\langle B,B\rangle_m}.
\end{equation}
Also retain $J_hF=\{h,F\}/2$, where
$F=W_1T_3-\bar t h$ and $\bar t=\langle h,W_1T_3\rangle_m/C$
are the previously defined linear rank-filtered field and its energy
projection. For direct evaluation in the same covariance primitives,
\begin{align}
 U_A&=U+(2p-m)(P_++Ch),\\
 W_1T_3&=\rr(2C_-k+P_-)+(3p-m)(2Ch+P_+)\nonumber\\
 &\hspace{12mm}-3\big(ah+(U_A+U_A^\dagger)/2\big).
\end{align}
All Gram entries can be evaluated without disorder sampling.
If a Hermitian quadratic polynomial has symmetric coefficients
$P=\sum_{ab}x_ax_bP_{ab}$, then
\begin{equation}
 \langle W_2P,W_2S\rangle_m
       =2\sum_{ab}v_av_b\tr_m(P_{ab}^\dagger S_{ab}).
 \label{eq:gram}
\end{equation}

At finite $N$, an additional conserved piece can be a random scalar
matrix. Remove it sample by sample with
$\pi X=X-\tr_m(X)I$. Project $\pi W_2\core$ onto the span of
$\pi B,\pi O,\pi J_hF$ and call the residual $\mathcal D$.
This removes a specified conserved subspace, not the full commutant.
Exact results at $w=1$ are
\begin{center}
\begin{tabular}{cccc}
\toprule
$(N,p,m)$ & $e_m$ & $c_E$ & $\langle\mathcal D,\mathcal D\rangle_m$\\
\midrule
$(3,1,2)$ & $38$ & $5$ & $0$\\
$(4,2,3)$ & $684$ & $1089/17$ & $0$\\
$(5,2,3)$ & $2268$ & $4596/41$ & $248945355/14903$\\
\bottomrule
\end{tabular}
\end{center}
The coefficient $c_E$ refers to the unmodified $B$ projection;
the last column includes $\pi$ and all three old directions.
For the last system their coefficients are, respectively,
\begin{equation}
 \left(\frac{2106991}{44709},\frac{1499753}{29806},
                 -\frac{77111}{89418}\right).
\end{equation}
The saved deterministic integer coupling vector gives
\begin{equation}
 \tr_m[h,\mathcal D]^\dagger[h,\mathcal D]
        =\frac{12991771132439678939}{2220994090}>0.
\end{equation}
Thus the extra finite-$N$ residual is genuinely dynamical in this
example. A positive finite residual alone does not determine its
double-scaled norm or spectral measure.

\newpage
\section{Analytic control and the next asymptotic calculation}
There is an all-$N$ one-body control. For $p=1$, let $J$ be the
fundamental $N\times N$ GUE matrix and $H=d\Gamma(J)$. The CAR give
$\mathcal L^2c=J^2c$, so
\begin{equation}
 A=d\Gamma(J^2),\quad [H,A]=0,\quad \core=d\Gamma(J^4).
\end{equation}
At $w=1$, $\EE KXK=\tr(X)I$ for the fundamental independent
matrix $K$, where this trace is unnormalized. The six contractions give
\begin{equation}
 Q_{\core}=3NA+2\tr(J)H+m\tr(J^2)I,\qquad
 e_m=m(2N^2+1).
\end{equation}
Every term commutes with $H$. This proves why a small one-body
example cannot establish a new dynamical mode, even when an
incomplete Gram subtraction leaves a positive norm.

The saved verification comprises three full finite-sector Gram
calculations (45, 666 and 5050 symmetric coefficient blocks),
807 independent quartic polarization directions, and 56 one-body
identities in 14 charge sectors. Polarization evaluates the original
commutator-defined core, independently of the charged covariance
expression. Gaussian coordinate variances are $1$ on diagonal
couplings and $1/2$ on each real or imaginary off-diagonal coordinate.
The Gram sums use rational arithmetic after checking dyadic matrix
coefficients; no Monte Carlo fit or charge truncation is used.

Restoring the physical variance rescales $W_2\core$ by $w^2$,
its squared norm by $w^4$, and the commutator witness by $w^5$.
For the intended double-scaled problem, first set $p^2/N\to\lambda>0$
at the physical $w$, then study small $\lambda$. The exact maps in
\eqref{eq:answer} should now be decomposed in their irreducible
input and output ranks. In particular, the $O(p)$ and $O(m)$
coefficients must be combined before discarding covariance differences.

The immediate target is the limiting norm and mixed energy moments
of the residual after restoring all fourth-chaos descendants. That
will decide whether it vanishes, joins existing filtered quadratic
channels, or supplies another direction. Next come its full two-energy
measure, the genuine $Z_2$ spectrum, and the covariance at two finite
memory parameters. Spatial reconstruction and the bulk mass dictionary
remain the subsequent physical objectives. No finite-$N$ channel
count is being identified here with a de Sitter particle spectrum.

\begin{thebibliography}{9}
\bibitem{Kota} V. K. B. Kota,
\emph{Bivariate moments of the two-point correlation function for embedded
Gaussian unitary ensemble with $k$-body interactions},
\href{https://arxiv.org/html/2208.11312}{arXiv:2208.11312v2},
Sec. IV, Eqs. (29)--(32).
\bibitem{TT} J. Tall and S. Tomsovic,
\emph{Double-scaled bosonic and fermionic embedded ensembles,
complex SYK, and the dual Hilbert space},
\href{https://arxiv.org/html/2604.14522v2}{arXiv:2604.14522v2},
Secs. 2.1 and 3.1.
\end{thebibliography}
\end{document}
